\section{Scope and retrospective audit design}
\label{sec:design}

We conduct a retrospective evidence audit of four campaigns from one
split-LLM research programme, followed by two checks on the scope of broader
security claims. The unit of analysis is a claim together with its supporting
records, rather than a token, a model, or a published paper. The campaigns
were selected because they supplied the synthesis being audited; they are
not a random sample of privacy evaluations. Consequently, their counts do
not estimate the prevalence of evaluation defects in the field. The
component results remain attributed to their original campaigns. Our
contribution is the reproducible reconciliation of those results with the
conclusions they admit.

This audit was specified after the historical outcomes were available. We
retain the distinction between a frozen exploratory protocol, a completed
experiment, a proposed confirmation, and the present secondary analysis.
Recomputing stored arithmetic does not make the underlying experiment
independently replicated or prospectively confirmatory. We do not retrain
victim models, collect new captures, fit stronger attackers, or infer
unrecorded outcomes. The evaluation-contract and control machinery of the
LLM-Attack Suite is relevant prior infrastructure\citep{suite2026}; this
paper audits cross-campaign claim admission rather than presenting that
machinery as a newly validated evaluation system.

\paragraph{Evidence-admission contract.}
For each claim we record the protocol and cohort; the adversary and its
access to secrets and remote state; the observable channels; the protected
target and its alignment with observations; the attack and training budget;
the score, baseline, and acceptance rule; the qualification binding; and the
independent unit and uncertainty calculation. A qualification binding
identifies the target, channel, defence condition, and implementation to
which positive and null controls apply. Passing a control elsewhere in the
programme does not qualify a new binding. A recorded channel is also not
automatically an evaluated attack surface.

We prefer the matching generated records and current admitted manifests
over narrative summaries. Where disposition changes, we retain the
supersession history and the later record's limitations rather than choosing
whichever account is more favourable. We distinguish quantities recomputed
from supplied data from results read through an existing audit or manuscript.
The accompanying source index and analysis bind each retained quantity to
its artifact. Those bindings establish traceability within the supplied
package; they do not certify the completeness of the original collection.

\paragraph{Five claim-specific outcomes.}
The following outcomes apply to a specified privacy claim and can coexist
across different endpoints of the same configuration:
\begin{description}
  \item[Breach.] An admissible, qualified measurement exceeds its declared
    privacy rule. The conclusion concerns that observer, target, attack, and
    finite protocol; it need not quantify all information disclosed.
  \item[Non-detection.] An admissible, qualified test does not exceed its
    rule. This is a result for the tested attack and budget. It is not an
    upper bound on leakage or a system-wide privacy guarantee.
  \item[Undetermined.] Missing coverage, incomplete qualification, or
    unavailable admissible observations prevents the required conclusion.
    An undetermined endpoint contributes no privacy pass.
  \item[Invalid.] A demonstrated defect breaks the connection between the
    measurement and the claimed estimand, as when next-token labels are
    used to claim current-token reconstruction resistance. The score may
    still describe a different, explicitly named quantity.
  \item[Out-of-scope.] The claim requires a different observer, secret
    access, target, or protocol from the one evaluated. An accurate
    measurement in the original setting cannot establish this claim.
\end{description}

Utility and remote-gain outcomes are recorded separately, and each
campaign's recorded aggregation rule is retained. In R1, for example, a
configuration can be rejected by a qualified privacy breach or an
independent utility failure while other privacy endpoints remain
undetermined. Missing qualification blocks a conjunctive pass; it does not
erase an admissible negative result. Conversely, failure to meet an
exploratory gain rule is not proof that the population gain is zero.

\paragraph{Limits on lifting a claim.}
We carry a conclusion to another setting only when the relevant protocol,
observer access, observation set, target, and measurement binding match, or
when a stated argument justifies the change. In particular, wire-observer
results cannot silently become compromised-host results, and empirical
non-detection cannot become an information-theoretic ceiling. Multiple
correlated rows do not increase the number of independent training seeds.
The four cases below illustrate these admission decisions; they do not
constitute an independent diagnostic-validation cohort or establish a
universal privacy decision procedure.
